%% kapa: exam=1 round=33 year=2022 period=2 form=A subject=accounting subject_ko=회계학 seq=10 number=50
%% source: Q-Net 감정평가사 기출문제 (cst00309) artlSeq=5209997 file=raw/1008/5209997/2022년도 제33회 감정평가사 1차시험 2교시(감정평가 관계법규, 회계학).pdf
%% conversion: kordoc 4.24.51 -> tools/gen1.py (auto); review=proofread-sonnet; flags=none
\begin{question}{50}
(주)감평은 재고상품에 대해 선입선출법을 적용하여 단위원가를 결정하며, 20×1년 기초상품은 ￦30,000(단위당 원가 ￦1,000), 당기상품매입액은 ￦84,000(단위당 원가 ￦1,200)이다. 기말상품의 감모손실과 평가손실에 관한 자료는 다음과 같다.
\par{\small\centering\begin{tabular}{|c|c|c|c|}\hline
장부수량 & 실제수량 & 단위당 예상판매가격 & 단위당 예상판매비용 \\ \hline
20개 & 16개 & ￦1,250 & ￦80 \\ \hline
\end{tabular}\par}
(주)감평이 기말 재고자산감모손실은 장부에 반영하였으나 재고자산평가손실을 반영하지 않았을 경우 옳은 것은?
\begin{choices}
  \item 20×1년 당기순이익 ￦1,000 과대
  \item 20×1년 기말재고자산 ￦600 과대
  \item 20×1년 기말자본총계 ￦480 과소
  \item 20×2년 기초재고자산 ￦600 과소
  \item 20×2년 당기순이익 ￦480 과소
\end{choices}
\end{question}
\answer{5}
